\documentclass[FileMain.tex]{subfiles}
\gdef\monhoc{Khoa học tự nhiên 6}
\gdef\ngaykt{Hè 2025 -- 2026}
\gdef\nh{2025 - 2026}
\gdef\thoigian{60}
\gdef\made{101}
\setcounter{section}{0}
%\tatloigiai
\hienthiloigiai
%\dongkeloigiai
\begin{document}
\section[Kiểm tra tổng hợp ôn hè Khoa học tự nhiên 6 - Mã đề \made]{Kiểm tra tổng hợp ôn hè}

%%%==============Phần trắc nghiệm nhiều lựa chọn==============%%%
\subsection{Bài tập trắc nghiệm nhiều lựa chọn}\textit{\large Thí sinh trả lời từ câu 1 đến câu 12. Mỗi câu thí sinh chỉ chọn một phương án}
\Opensolutionfile{ansex}[Ans/LGEX-DE_TON_HE_HOA6_MADE101]
\Opensolutionfile{ans}[Ans/Ans-DE_TON_HE_HOA6_MADE101]

%%%=============EX_1=============%%%
\begin{ex}%[6K2N2-1]
Khí nào chiếm tỉ lệ lớn nhất trong thành phần không khí?
\choice
{Oxygen ($O_2$)}
{\True Nitrogen ($N_2$)}
{Carbon dioxide ($CO_2$)}
{Hơi nước ($H_2O$)}
\loigiai{
 Thành phần không khí gồm: Nitrogen ($N_2$) chiếm khoảng $78\%$, Oxygen ($O_2$) chiếm khoảng $21\%$, còn lại là $CO_2$, hơi nước và các khí khác. Vậy Nitrogen chiếm tỉ lệ lớn nhất.
}
\end{ex}

%%%=============EX_2=============%%%
\begin{ex}%[6K2N2-3]
Tính chất nào sau đây là của không khí?
\choice
{Có màu trắng đục}
{Có hình dạng cố định}
{\True Có thể bị nén hoặc dãn nở}
{Có vị ngọt đặc trưng}
\loigiai{
 Không khí trong suốt, không màu, không mùi, không vị; không có hình dạng và thể tích nhất định; có thể bị nén lại hoặc dãn ra. Đây là tính chất đặc trưng của chất khí.
}
\end{ex}

%%%=============EX_3=============%%%
\begin{ex}%[6K1N2-2]
Hiện tượng nước chuyển từ thể lỏng sang thể khí được gọi là gì?
\choice
{Nóng chảy}
{Đông đặc}
{Ngưng tụ}
{\True Bay hơi}
\loigiai{
 Các hiện tượng chuyển thể: Lỏng $\to$ Khí: Bay hơi; Khí $\to$ Lỏng: Ngưng tụ; Rắn $\to$ Lỏng: Nóng chảy; Lỏng $\to$ Rắn: Đông đặc.
}
\end{ex}

%%%=============EX_4=============%%%
\begin{ex}%[6K1N2-1]
Nến đang cháy, phần nến lỏng xung quanh bấc nguội dần và đóng cứng lại. Hiện tượng này là sự
\choice
{bay hơi}
{ngưng tụ}
{\True đông đặc}
{nóng chảy}
\loigiai{
 Nến lỏng nguội dần và trở về trạng thái rắn là sự đông đặc (chất chuyển từ thể lỏng sang thể rắn).
}
\end{ex}

%%%=============EX_5=============%%%
\begin{ex}%[6K1N1-2]
Chất ở trạng thái khí có đặc điểm nào sau đây?
\choice
{Có hình dạng và thể tích xác định}
{Có hình dạng không xác định, thể tích xác định}
{\True Không có hình dạng và thể tích xác định, chiếm toàn bộ không gian vật chứa}
{Có hình dạng xác định, thể tích không xác định}
\loigiai{
 Chất khí: hình dạng không xác định, chiếm toàn bộ không gian vật chứa (thể tích không xác định). Chất lỏng: hình dạng không xác định, thể tích xác định. Chất rắn: hình dạng và thể tích đều xác định.
}
\end{ex}

%%%=============EX_6=============%%%
\begin{ex}%[6K4N1-1]
Hỗn hợp nào sau đây là dung dịch?
\choice
{Gạo trộn lẫn với sỏi}
{Cát và nước khuấy đều}
{\True Muối hòa tan trong nước sau khi khuấy đều}
{Dầu ăn và nước khuấy đều}
\loigiai{
 Dung dịch là hỗn hợp chất lỏng với chất rắn (hoặc chất lỏng) hòa tan và phân bố đều vào nhau. Muối hòa tan hoàn toàn trong nước tạo thành dung dịch đồng nhất. Cát và dầu ăn không tan trong nước nên không tạo thành dung dịch.
}
\end{ex}

%%%=============EX_7=============%%%
\begin{ex}%[6K4N1-2]
Đặc điểm nào phân biệt hỗn hợp với dung dịch?
\choice
{\True Trong hỗn hợp, các chất giữ nguyên tính chất; trong dung dịch, tính chất các chất bị biến đổi}
{Hỗn hợp luôn trong suốt, dung dịch thì đục}
{Hỗn hợp có màu sắc, dung dịch không có màu}
{Hỗn hợp chứa nhiều chất hơn dung dịch}
\loigiai{
 Trong hỗn hợp, các chất trộn lẫn nhưng vẫn giữ nguyên tính chất. Trong dung dịch, các chất hòa tan và phân bố đều vào nhau nên tính chất bị biến đổi.
}
\end{ex}

%%%=============EX_8=============%%%
\begin{ex}%[6K4N3-2]
Để tách muối ra khỏi dung dịch nước muối, người ta dùng phương pháp nào?
\choice
{Lọc qua vải}
{Dùng nam châm}
{\True Đun nóng để bay hơi nước (cô cạn)}
{Để yên cho muối tự lắng xuống}
\loigiai{
 Vì muối không bay hơi còn nước bay hơi khi đun nóng, ta dùng phương pháp cô cạn (đun nóng dung dịch cho nước bay hơi hết) để thu được muối kết tinh.
}
\end{ex}

%%%=============EX_9=============%%%
\begin{ex}%[6K1N1-4]
Hiện tượng nào sau đây là sự biến đổi hóa học?
\choice
{Nước đá tan chảy thành nước lỏng}
{Xé giấy thành nhiều mảnh nhỏ}
{Hòa tan đường vào nước}
{\True Đinh sắt để lâu ngoài không khí ẩm bị gỉ, chuyển sang màu nâu đỏ}
\loigiai{
 Biến đổi hóa học là biến đổi có sự tạo thành chất mới. Đinh sắt bị gỉ tạo ra chất mới (gỉ sắt) có màu nâu đỏ và tính chất khác sắt ban đầu. Các hiện tượng còn lại chỉ thay đổi hình dạng hoặc trạng thái, không tạo chất mới.
}
\end{ex}

%%%=============EX_10=============%%%
\begin{ex}%[6K1N1-6]
Dấu hiệu nào giúp nhận biết đã xảy ra biến đổi hóa học?
\choice
{Chất thay đổi hình dạng}
{Chất chuyển từ trạng thái lỏng sang rắn}
{\True Xuất hiện chất mới với màu sắc, mùi vị hoặc tính tan khác}
{Chất bị chia nhỏ ra}
\loigiai{
 Biến đổi hóa học tạo ra chất mới có tính chất khác chất ban đầu, biểu hiện qua: thay đổi màu sắc, xuất hiện mùi lạ, thay đổi tính tan,\dots Các dấu hiệu còn lại là biến đổi vật lý.
}
\end{ex}

%%%=============EX_11=============%%%
\begin{ex}%[6K8N1-5]
Khi đun nóng ống nghiệm trên ngọn lửa đèn cồn, cần lưu ý điều gì?
\choice
{Để miệng ống nghiệm hướng vào người ngồi bên cạnh}
{\True Dùng kẹp gỗ giữ ống nghiệm và hướng miệng ống về phía không có người}
{Cầm trực tiếp ống nghiệm bằng tay}
{Đặt ống nghiệm thẳng đứng, miệng hướng lên trên}
\loigiai{
 Quy tắc an toàn: Khi đun nóng ống nghiệm phải dùng kẹp gỗ giữ và hướng miệng ống về phía không có người, để tránh bị thương nếu chất lỏng sôi và bắn ra.
}
\end{ex}

%%%=============EX_12=============%%%
\begin{ex}%[6K8N1-6]
Dụng cụ nào dùng để đo chính xác thể tích chất lỏng trong phòng thí nghiệm?
\choice
{Bình cầu}
{Cốc thủy tinh}
{\True Ống đong có vạch chia}
{Đũa thủy tinh}
\loigiai{
 Ống đong có vạch chia được dùng để đo thể tích chất lỏng chính xác. Cốc thủy tinh và bình cầu chỉ dùng để đựng/đun chất lỏng, không có độ chính xác cao. Đũa thủy tinh dùng để khuấy.
}
\end{ex}

\Closesolutionfile{ans}
\Closesolutionfile{ansex}

%%%==============Phần trắc nghiệm đúng sai==============%%%
\subsection{Trắc nghiệm đúng sai}\textit{\large Thí sinh trả lời từ câu 1 đến câu 4. Trong mỗi ý a), b), c), d) ở mỗi câu thí sinh chọn đúng hoặc sai}
\Opensolutionfile{ansex}[Ans/LGTF-DE_TON_HE_HOA6_MADE101]
\Opensolutionfile{ansbook}[Ansbook/AnsTF-DE_TON_HE_HOA6_MADE101]
\Opensolutionfile{ans}[Ans/Tempt-DE_TON_HE_HOA6_MADE101]
\setcounter{ex}{0}

%%%=============TF_1=============%%%
\begin{ex}%[6K1H2-5]
Nước là chất quen thuộc trong đời sống. Xét các tính chất và hiện tượng liên quan đến nước, đánh giá tính đúng/sai của các phát biểu sau:
\choiceTF
{\True Nước tinh khiết sôi ở $100^\circ C$ và đông đặc ở $0^\circ C$ ở điều kiện bình thường}
{\True Hiện tượng sương mù buổi sáng là do hơi nước trong không khí gặp lạnh ngưng tụ thành các giọt nước nhỏ}
{Ở vùng núi cao, nước sôi ở nhiệt độ cao hơn $100^\circ C$ vì không khí trên núi lạnh hơn}
{Nước không có hình dạng cố định và cũng không có thể tích xác định}
\loigiai{
 \begin{itemchoice}[T1,T2,F3,F4]
 \itemch Ở điều kiện tiêu chuẩn (áp suất khí quyển 1 atm), nước tinh khiết sôi ở $100^\circ C$ và đông đặc ở $0^\circ C$. Phát biểu đúng.
 \itemch Sương mù buổi sáng là hơi nước trong không khí ấm gặp không khí lạnh gần mặt đất, ngưng tụ thành giọt nước nhỏ lơ lửng. Đây là sự ngưng tụ. Phát biểu đúng.
 \itemch Ở vùng núi cao, áp suất khí quyển thấp hơn nên nhiệt độ sôi của nước \textit{thấp hơn} $100^\circ C$, không phải cao hơn. Phát biểu sai.
 \itemch Nước ở trạng thái lỏng không có hình dạng cố định nhưng \textit{có thể tích xác định} (tính chất của chất lỏng). Phát biểu sai.
 \end{itemchoice}
}
\end{ex}

%%%=============TF_2=============%%%
\begin{ex}%[6K4H1-3]
Một học sinh thực hiện các thí nghiệm với các hỗn hợp khác nhau. Đánh giá tính đúng/sai của các nhận xét sau:
\choiceTF
{Khi trộn cát và nước khuấy đều, ta thu được một dung dịch trong suốt}
{\True Nước đường là dung dịch vì đường hòa tan hoàn toàn và phân bố đều trong nước}
{\True Để tách mạt sắt ra khỏi hỗn hợp mạt sắt và cát, người ta có thể dùng nam châm vì mạt sắt bị nam châm hút}
{Trong dung dịch, tính chất của các chất thành phần vẫn được giữ nguyên như ban đầu}
\loigiai{
 \begin{itemchoice}[F1,T2,T3,F4]
 \itemch Cát không hòa tan trong nước, chỉ lơ lửng rồi lắng xuống. Đây là hỗn hợp không đồng nhất (huyền phù), không phải dung dịch. Phát biểu sai.
 \itemch Đường hòa tan hoàn toàn và phân bố đều trong nước tạo thành dung dịch đồng nhất, trong suốt. Phát biểu đúng.
 \itemch Sắt có tính từ (bị nam châm hút), cát không bị hút. Dùng nam châm có thể tách mạt sắt ra khỏi cát. Phát biểu đúng.
 \itemch Trong dung dịch, các chất hòa tan vào nhau khiến tính chất bị biến đổi (ví dụ: nước muối mặn hơn nước nguyên chất). Phát biểu sai.
 \end{itemchoice}
}
\end{ex}

%%%=============TF_3=============%%%
\begin{ex}%[6K1H1-5]
Xét các hiện tượng trong đời sống, đánh giá tính đúng/sai của các phát biểu sau về biến đổi chất:
\choiceTF
{\True Cơm bị cháy khét khi nấu quá lửa là biến đổi hóa học vì tạo ra chất mới có màu đen và mùi khét}
{Nước đá tan chảy thành nước lỏng là biến đổi hóa học vì chất đã thay đổi trạng thái}
{\True Bảo quản thực phẩm trong tủ lạnh giúp làm chậm sự biến đổi hóa học khiến thức ăn bị ôi thiu}
{\True Vắt chanh vào nước rau muống làm nước chuyển từ màu xanh sang đỏ nhạt là biến đổi hóa học}
\loigiai{
 \begin{itemchoice}[T1,F2,T3,T4]
 \itemch Cơm cháy tạo ra chất mới (carbon và các sản phẩm cháy) có màu đen, mùi khét, khác hoàn toàn cơm ban đầu. Đây là biến đổi hóa học. Phát biểu đúng.
 \itemch Nước đá $\to$ nước lỏng chỉ là sự thay đổi trạng thái (nóng chảy), không tạo ra chất mới. Đây là biến đổi vật lý. Phát biểu sai.
 \itemch Nhiệt độ thấp trong tủ lạnh làm chậm quá trình biến đổi hóa học của vi khuẩn và enzyme trong thực phẩm. Phát biểu đúng.
 \itemch Acid trong chanh phản ứng với chất màu trong rau muống tạo thành sản phẩm mới có màu đỏ nhạt. Đây là biến đổi hóa học. Phát biểu đúng.
 \end{itemchoice}
}
\end{ex}

%%%=============TF_4=============%%%
\begin{ex}%[6K8H1-5]
Một học sinh đang thực hành trong phòng thí nghiệm. Đánh giá tính đúng/sai của các hành động và quy tắc dưới đây:
\choiceTF
{\True Khi hóa chất bắn vào mắt, việc đầu tiên cần làm là rửa mắt ngay bằng nước sạch nhiều lần rồi báo giáo viên}
{\True Để tắt đèn cồn an toàn, cần dùng nắp đậy kín ngọn lửa, không dùng miệng thổi}
{Khi làm thí nghiệm, có thể ăn nhẹ trong phòng thí nghiệm nếu đang nghỉ giải lao}
{Hóa chất thừa sau thí nghiệm có thể đổ trực tiếp vào bồn rửa mà không cần xử lý}
\loigiai{
 \begin{itemchoice}[T1,T2,F3,F4]
 \itemch Khi hóa chất bắn vào mắt phải rửa ngay bằng nước sạch liên tục để pha loãng và loại bỏ hóa chất, sau đó báo giáo viên. Phát biểu đúng.
 \itemch Thổi miệng vào đèn cồn có thể làm cồn bắn ra gây cháy. Đúng cách là dùng nắp đậy để ngọn lửa tắt do thiếu oxygen. Phát biểu đúng.
 \itemch Tuyệt đối không ăn uống trong phòng thí nghiệm vì hóa chất có thể bám vào thức ăn gây ngộ độc. Phát biểu sai.
 \itemch Hóa chất thừa phải thu gom và xử lý theo hướng dẫn của giáo viên, không đổ trực tiếp vào bồn rửa vì có thể gây ô nhiễm môi trường và hư hỏng đường ống. Phát biểu sai.
 \end{itemchoice}
}
\end{ex}

\Closesolutionfile{ans}
\Closesolutionfile{ansbook}
\Closesolutionfile{ansex}

%%==============Phần bài tập trả lời ngắn==============%%%
\subsection{Bài tập trả lời ngắn}\textit{\large Thí sinh trả lời từ câu 1 đến câu 3}
\Opensolutionfile{ansex}[Ans/LGSA-DE_TON_HE_HOA6_MADE101]
\Opensolutionfile{ansexh}[Ans/AnsSA-DE_TON_HE_HOA6_MADE101]
\setcounter{ex}{0}

%%%=============SA_1=============%%%
\begin{ex}%[6K2N2-1]
Oxygen chiếm bao nhiêu phần trăm thể tích không khí? (Điền số nguyên)
\shortans{$21$}
\loigiai{
 Thành phần không khí theo thể tích: Nitrogen ($N_2$) chiếm khoảng $78\%$, Oxygen ($O_2$) chiếm khoảng $21\%$, còn lại khoảng $1\%$ là $CO_2$, hơi nước và các khí khác.
}
\end{ex}

%%%=============SA_2=============%%%
\begin{ex}%[6K1N2-3]
Ở điều kiện bình thường (áp suất khí quyển), nước tinh khiết sôi ở nhiệt độ bao nhiêu độ C?
\shortans{$100$}
\loigiai{
 Ở áp suất khí quyển tiêu chuẩn ($1$ atm), nước tinh khiết sôi ở $100^\circ C$. Đây là nhiệt độ sôi tiêu chuẩn của nước.
}
\end{ex}

%%%=============SA_3=============%%%
\begin{ex}%[6K1N2-1]
Nhiệt độ đông đặc của nước tinh khiết ở điều kiện bình thường là bao nhiêu độ C?
\shortans{$0$}
\loigiai{
 Ở áp suất khí quyển tiêu chuẩn, nước tinh khiết đông đặc (chuyển từ thể lỏng sang thể rắn) ở $0^\circ C$.
}
\end{ex}

\Closesolutionfile{ansexh}
\Closesolutionfile{ansex}

%%%==============Phần bài tập tự luận==============%%%
\subsection{Bài tập tự luận}\textit{\large Thí sinh trả lời từ bài 1 đến bài 2}
\Opensolutionfile{ansbth}[Ans/LGBT-DE_TON_HE_HOA6_MADE101]
\Opensolutionfile{ansbt}[Ans/AnsBT-DE_TON_HE_HOA6_MADE101]

%%%=============BT_1=============%%%
\begin{bt}%[6K1V2-4]
Quan sát sơ đồ vòng tuần hoàn của nước trong tự nhiên:
\begin{center}
\textit{Nước biển $\xrightarrow[$\text{nhiệt mặt trời}$][$\text{bay hơi$]}$ Hơi nước $\xrightarrow[$\text{gặp lạnh}$][$\text{ngưng tụ$]}$ Mây $\xrightarrow[$$]$ Mưa $\xrightarrow[$$]$ Sông/suối $\xrightarrow[$$]$ Biển}
\end{center}
\begin{enumerate}[(a)]
 \item Kể tên các sự chuyển thể của nước diễn ra trong vòng tuần hoàn trên.
 \item Giải thích: Tại sao ở vùng núi cao, nước sôi ở nhiệt độ thấp hơn $100^\circ C$?
\end{enumerate}
\loigiai{
\textbf{(a) Các sự chuyển thể trong vòng tuần hoàn của nước:}
\begin{itemize}
 \item \textbf{Bay hơi}: Nước ở biển, sông, hồ chuyển từ thể lỏng sang thể khí dưới tác động của nhiệt mặt trời.
 \item \textbf{Ngưng tụ}: Hơi nước bốc lên cao gặp lạnh ngưng tụ thành các giọt nước nhỏ tạo thành mây (khí $\to$ lỏng).
 \item \textbf{Mưa}: Các giọt nước trong mây lớn dần và rơi xuống dưới dạng thể lỏng.
 \item \textbf{Đông đặc} (ở vùng núi cao/cực): Nước có thể đông thành tuyết, băng (lỏng $\to$ rắn).
 \item \textbf{Nóng chảy}: Tuyết/băng tan chảy thành nước lỏng (rắn $\to$ lỏng).
\end{itemize}

\textbf{(b) Giải thích tại sao nước sôi ở nhiệt độ thấp hơn $100^\circ C$ trên núi cao:}

Trên núi cao, áp suất khí quyển thấp hơn so với mực nước biển. Nhiệt độ sôi của chất lỏng phụ thuộc vào áp suất tác dụng lên bề mặt: áp suất càng thấp thì nhiệt độ sôi càng thấp. Do đó, trên núi cao, nước sôi ở nhiệt độ thấp hơn $100^\circ C$.
}
\end{bt}

%%%=============BT_2=============%%%
\begin{bt}%[6K1V1-6]
Cho các hiện tượng sau:
\begin{enumerate}[(1)]
 \item Nước đá tan chảy thành nước lỏng khi để ra ngoài trời nắng.
 \item Đinh sắt để lâu ngoài không khí ẩm bị gỉ, chuyển sang màu nâu đỏ và trở nên giòn, dễ gãy.
 \item Đường cát hòa tan vào nước khi khuấy đều.
 \item Đốt cháy một tờ giấy, giấy biến thành tro và khói.
\end{enumerate}
\begin{enumerate}[(a)]
 \item Phân loại các hiện tượng trên thành biến đổi vật lý và biến đổi hóa học.
 \item Nêu dấu hiệu để em nhận biết biến đổi hóa học ở hiện tượng (2) và (4).
\end{enumerate}
\loigiai{
\textbf{(a) Phân loại:}
\begin{itemize}
 \item \textbf{Biến đổi vật lý} (không tạo ra chất mới):
 \begin{itemize}
 \item Hiện tượng (1): Nước đá tan --- chỉ thay đổi trạng thái (rắn $\to$ lỏng), vẫn là nước $H_2O$.
 \item Hiện tượng (3): Đường hòa tan --- đường vẫn là đường, chỉ phân tán đều trong nước.
 \end{itemize}
 \item \textbf{Biến đổi hóa học} (tạo ra chất mới):
 \begin{itemize}
 \item Hiện tượng (2): Đinh sắt bị gỉ --- tạo ra gỉ sắt là chất mới.
 \item Hiện tượng (4): Đốt giấy --- giấy biến thành tro và khí $CO_2$, là các chất mới.
 \end{itemize}
\end{itemize}

\textbf{(b) Dấu hiệu nhận biết biến đổi hóa học:}
\begin{itemize}
 \item \textbf{Hiện tượng (2) --- Đinh sắt gỉ:}
 \begin{itemize}
 \item Thay đổi \textbf{màu sắc}: từ màu xám sáng của sắt sang màu nâu đỏ của gỉ sắt.
 \item Thay đổi \textbf{tính chất cơ học}: từ bền chắc sang giòn, dễ gãy. Đây là dấu hiệu chất mới được tạo thành.
 \end{itemize}
 \item \textbf{Hiện tượng (4) --- Đốt giấy:}
 \begin{itemize}
 \item Xuất hiện \textbf{màu sắc mới}: giấy trắng $\to$ tro đen.
 \item Xuất hiện \textbf{mùi}: khói có mùi khét đặc trưng.
 \item Tạo ra sản phẩm hoàn toàn khác chất ban đầu: tro và khí --- là những chất mới.
 \end{itemize}
\end{itemize}
}
\end{bt}

\Closesolutionfile{ansbt}
\Closesolutionfile{ansbth}

\begin{center}
 \rule[4pt]{2cm}{1pt}\,\large\bfseries Hết\,\rule[4pt]{2cm}{1pt}
\end{center}
\label{x}
\end{document}
